\subsection{习题}
\label{sec:ch14-exercises}
\setcounter{exerciseitem}{0}

\begin{MBExercise}
\label{exr:ch14-01}
设 $A_i$ 和 $B_i$ 是上述两对 Banach 空间, 且连续嵌入 $B_i\subset A_i$($i=0,1$). 证明 $B_{\theta,p}\subset A_{\theta,p}$. (提示: 设 $C$ 使 $\|v\|_{A_i}\leq C\|v\|_{B_i}$, 并证明 $K_A(t,u)\leq CK_B(t,u)$. )
\end{MBExercise}

\begin{MBExercise}
\label{exr:ch14-02}
证明包含算子~\eqref{eq:ch14-domain-restriction-interpolation-inclusion} 的范数为 $1$.
\end{MBExercise}

\begin{MBExercise}
\label{exr:ch14-03}
证明~\eqref{eq:ch14-k-functional-interpolation-bound}. (提示: 注意 $\min\{1,s/t\}K(t,u)\leq K(s,u)$, 并对该不等式积分. )
\end{MBExercise}

\begin{MBExercise}
\label{exr:ch14-04}
已知 $B_1\subset B_0$, 证明 $B_1\subset B_{\theta,p}\subset B_0$, 且所有包含映射都连续. (提示: 假设 $\|v\|_{B_0}\leq C\|v\|_{B_1}$, 并证明 $\|u\|_{B_0}\leq K(C,u)$, 再用习题~\ref{exr:ch14-03} 验证第二个包含关系. 为证明第一个包含关系, 利用~\eqref{eq:ch14-real-interpolation-norm} 之前给出的 $v$ 的选择, 证明 $K(t,u)\leq\min\{C,t\}\|u\|_{B_1}$. )
\end{MBExercise}

\begin{MBExercise}
\label{exr:ch14-05}
证明~\eqref{eq:ch14-k-functional-little-o}. (提示: 注意 $t^{-\theta}K(t,u)\in L^p(0,\infty;dt/t)$, 因而当 $t\to0$ 时它必须趋于零. )
\end{MBExercise}

\MathBookSourcePage{389}
\begin{MBExercise}
\label{exr:ch14-06}
当 $s$ 为整数时, 证明~\eqref{eq:ch14-fourier-sobolev-norm} 是 $H^s(\Omega)$ 的一个等价范数. (提示: 回顾 $\widehat{D^\alpha u}=(i\xi)^\alpha\widehat u$, 且 Fourier 变换是 $L^2$ 中的同构. 注意若 $0<r<s$, 则 $|\xi|^r\leq1+|\xi|^s$. )
\end{MBExercise}

\begin{MBExercise}
\label{exr:ch14-07}
当 $s$ 为分数时, 证明~\eqref{eq:ch14-fourier-sobolev-norm} 与~\eqref{eq:ch14-intrinsic-fractional-sobolev-norm} 等价. 提示: 使用 Plancherel 定理和 Fubini 定理写成
\[
\MathBookFormulaID{formula-ch14-exercise-fractional-fourier-identity}
\begin{aligned}
\int_{\mathbb R^n}\int_{\mathbb R^n}
\frac{|u(x+r)-u(x)|^2}{|r|^{d+2s}}\,dx\,dr
&=\int_{\mathbb R^n}\int_{\mathbb R^n}
|\widehat u(\xi)(e^{ir\cdot\xi}-1)|^2\,d\xi\,
\frac{dr}{|r|^{d+2s}},
\end{aligned}
\]
并证明对 $\xi\in\mathbb R^n$ 和 $0<s<1$,
\[
\MathBookFormulaID{formula-ch14-exercise-fractional-fourier-bound}
\int_{\mathbb R^n}\frac{|e^{ir\cdot\xi}-1|^2}{|r|^{d+2s}}\,dr
\leq C(1+|\xi|^2)^s.
\]
\end{MBExercise}

\begin{MBExercise}
\label{exr:ch14-08}
设 $\Omega$ 为区间 $[-1,1]$, 并令 $f$ 为 Heavyside 函数: 当 $x<0$ 时 $f\equiv0$, 当 $x>0$ 时 $f\equiv1$. 证明 $f\in[L^2(\Omega),H^1(\Omega)]_{1/2,\infty}$. 提示: 对每个 $1>t>0$, 令 $v$ 为连续分片线性函数, 满足当 $x<-t$ 时 $v\equiv0$, 当 $x>t$ 时 $v\equiv1$, 且在 $[-t,t]$ 上为线性函数. 用此函数证明
\[
\MathBookFormulaID{formula-ch14-exercise-heavyside-k-bound}
K(t,f)\leq\|f-v\|_{L^2(\Omega)}+t\|v\|_{H^1(\Omega)}\leq Ct^{1/2}.
\]
\end{MBExercise}
